\begin{textsample}{POS Dim 1 – 2019 – Score 21.00 – 2019/t001498.txt}  \label{ex:f1_pos_057}
My mom has always reminded me that I have \textbf{the} same proportions as a LEGO man.
And she does actually have a point.
LEGO is a company that has succeeded in making everybody believe that LEGO is from their home country.
But it ' s not, it ' s from my home country.
So you can imagine my excitement when \textbf{the} LEGO family called me and asked us to work with them to design \textbf{the} Home of \textbf{the} Brick.
This is \textbf{the} architectural model — we built it out of LEGO, obviously.
This is \textbf{the} final result.
And what we tried to do was to design a building that would be as interactive and as engaging and as playful as LEGO is itself, with these kind of interconnected playgrounds on \textbf{the} roofscape.
You can enter a square on \textbf{the} ground where \textbf{the} citizens of Billund can roam around freely without a \textbf{ticket}.
And it ' s probably one of \textbf{the} only museums in \textbf{the} world where you ' re allowed to touch all \textbf{the} artifacts.
But \textbf{the} Danish word for design is"formgivning,"which literally means to give form to that which has not yet been given form.
In other words, to give form to \textbf{the} future.
And what I love about LEGO is that LEGO is not a toy.
It ' s a tool that empowers \textbf{the} child to build his or her own world, and then to inhabit that world through play and to invite her friends to join her in cohabiting and cocreating that world.
And that is exactly what formgivning is.
As human beings, we have \textbf{the} power to give form to our future.
Inspired by LEGO, we ' ve built a social housing project in Copenhagen, where we \textbf{stacked} blocks of wood next to each other.
Between them, they leave spaces with extra ceiling heights and balconies.
And by gently wiggling \textbf{the} blocks, we can actually create curves or any \textbf{organic} form, adapting to any urban context.
Because adaptability is probably one of \textbf{the} strongest drivers of architecture.
Another example is here in Vancouver.
We were asked to look at \textbf{the} site where Granville bridge triforks as it touches downtown.
And we started, like, mapping \textbf{the} different constraints.
There ' s like a 100-foot setback from \textbf{the} bridge because \textbf{the} city want to make sure that no one looks into \textbf{the} traffic on \textbf{the} bridge.
There ' s a park where we can ' t cast any shadows.
So finally, we ' re left with a tiny triangular footprint, almost too small to build.
But then we thought, like, what if \textbf{the} 100-foot minimum distance is really about minimum distance — once we get 100 feet up in \textbf{the} air, we can grow \textbf{the} building back out.
And so we did.
When you drive over \textbf{the} bridge, it ' s as if someone is pulling a curtain aback, welcoming you to Vancouver.
Or a like a weed growing through \textbf{the} cracks in \textbf{the} pavement and blossoming as it gets light and air.
Underneath \textbf{the} bridge, we ' ve worked with Rodney Graham and a handful of Vancouver artists, to create what we called \textbf{the} Sistine Chapel of street art, an art gallery turned upside down, that tries to turn \textbf{the} negative impact of \textbf{the} bridge into a positive.
So even if it looks like this kind of surreal architecture, it ' s highly adapted to its surroundings.
So if a bridge can become a museum, a museum can also serve as a bridge.
In Norway, we are building a museum that spans across a river and allows people to sort of journey through \textbf{the} exhibitions as they cross from one side of a sculpture park to \textbf{the} other.
An architecture sort of adapted to its landscape.
In China, we built a headquarters for an energy company and we designed \textbf{the} facade like an Issey Miyake fabric.
It ' s rippled, so that facing \textbf{the} predominant direction of \textbf{the} sun, it ' s all opaque ; facing away from \textbf{the} sun, it ' s all glass.
On average, it sort of transitions from solid to clear.
And this very simple idea without any moving parts or any sort of technology, purely because of \textbf{the} geometry of \textbf{the} facade, reduces \textbf{the} energy \textbf{consumption} on cooling by 30 percent.
So you can say what makes \textbf{the} building look elegant is also what makes it perform elegantly.
It ' s an architecture that is adapted to its climate.
You can also adapt one culture to another, like in Manhattan, we took \textbf{the} Copenhagen courtyard building with a social space where people can hang out in this kind of oasis in \textbf{the} middle of a city, and we combined it with \textbf{the} density and \textbf{the} verticality of an American skyscraper, creating what we ' ve called a"courtscraper."From New York to Copenhagen.
On \textbf{the} waterfront of Copenhagen, we are right now finishing this waste-to-energy power plant.
It ' s going to be \textbf{the} cleanest waste-to-energy power plant in \textbf{the} world, there are no toxins coming out of \textbf{the} chimney.
An amazing marvel of engineering that is completely invisible.
So we thought, how can we express this?
And in Copenhagen we have snow, as you can see, but we have absolutely no mountains.
We have to go six hours by bus to get to Sweden, to get alpine skiing.
So we thought, let ' s put an alpine ski \textbf{slope} on \textbf{the} roof of \textbf{the} power plant.
So this is \textbf{the} first test run we did a few months ago.
And what I like about this is that it also show you \textbf{the} sort of world-changing power of formgivning.
I have a five-month-old son, and he ' s going to grow up in a world not knowing that there was ever a time when you couldn ' t ski on \textbf{the} roof of \textbf{the} power plant.
So imagine for him and his generation, that ' s their baseline.
Imagine how far they can leap, what kind of wild ideas they can put forward for their future.
So right in front of it, we ' re building our smallest project.
It ' s basically nine containers that we have \textbf{stacked} in a shipyard in Poland, then we ' ve schlepped it across \textbf{the} Baltic \textbf{sea} and docked it in \textbf{the} port of Copenhagen, where it is now \textbf{the} home of 12 students.
Each student has a view to \textbf{the} water, they can jump out \textbf{the} window into \textbf{the} clean port of Copenhagen, and they can get back in.
All of \textbf{the} heat comes from \textbf{the} thermal mass of \textbf{the} \textbf{sea}, all \textbf{the} power comes from \textbf{the} sun.
This is \textbf{the} first 12 units in Copenhagen, another 60 on their way, another 200 are going to Gothenburg, and we ' re speaking with \textbf{the} Paris Olympics to put a small floating village on \textbf{the} Seine.
So very much this kind of, almost like nomadic, impermanent architecture.
And \textbf{the} waterfronts of our cities are experiencing a lot of change.
Economic change, industrial change and climate change.
This is Manhattan before Hurricane Sandy, and this is Manhattan after Sandy.
We got invited by \textbf{the} city of New York to look if we could make \textbf{the} necessary flood protection for Manhattan without building a seawall that would segregate \textbf{the} life of \textbf{the} city from \textbf{the} water around it.
And we got inspired by \textbf{the} High Line.
You probably know \textbf{the} High Line — it ' s this amazing new park in New York.
It ' s basically decommissioned train tracks that now have become one of \textbf{the} most popular promenades in \textbf{the} city.
So we thought, could we design \textbf{the} necessary flood protection for Manhattan so we don ' t have to wait until we shut it down before it gets nice?
So we sat down with \textbf{the} citizens living along \textbf{the} waterfront of New York, and we worked with them to try to design \textbf{the} necessary flood protection in such a way that it only makes their waterfront more accessible and more enjoyable.
Underneath \textbf{the} FDR, we are putting, like, pavilions with \textbf{pocket} walls that can \textbf{slide} out and protect from \textbf{the} water.
We are creating little stepped terraces that are going to make \textbf{the} underside more enjoyable, but also protect from flooding.
Further north in \textbf{the} East River Park, we are creating rolling hills that protect \textbf{the} park from \textbf{the} noise of \textbf{the} \textbf{highway}, but in turn also become \textbf{the} necessary flood protection that can stop \textbf{the} waves during an incoming storm surge.
So in a way, this project that we have called \textbf{the} Dryline, it ' s essentially \textbf{the} High Line — \textbf{The} High Line that ' s going to keep Manhattan dry.
It ' s scheduled to break ground on \textbf{the} first East River portion at \textbf{the} end of this year.
But it has essentially been codesigned with \textbf{the} citizens of Lower Manhattan to take all of \textbf{the} necessary infrastructure for resilience and give it positive social and environmental side effects.
So, New York is not alone in facing this situation.
In fact, by 2050, 90 percent of \textbf{the} major cities in \textbf{the} world are going to be dealing with rising \textbf{seas}.
In Hamburg, they ' ve created a whole neighborhood where \textbf{the} bottom floors are designed to withstand \textbf{the} inevitable flood.
In Sweden, they ' ve designed a city where all of \textbf{the} parks are \textbf{wet} gardens, designed to deal with storm water and waste water.
So we thought, could we perhaps — Actually, today, three million people are already permanently living on \textbf{the} \textbf{sea}.
So we thought, could we actually imagine a floating city designed to incorporate all of \textbf{the} Sustainable Development Goals of \textbf{the} United Nations into a whole new human-made ecosystem.
And of course, we have to design it so it can produce its own power, harvesting \textbf{the} thermal mass of \textbf{the} oceans, \textbf{the} force of \textbf{the} tides, of \textbf{the} currents, of \textbf{the} waves, \textbf{the} power of \textbf{the} wind, \textbf{the} heat and \textbf{the} energy of \textbf{the} sun.
Also, we are going to collect all of \textbf{the} rain water that drops on this man-made archipelago and deal with it organically and mechanically and store it and clean it.
We have to grow all of our food locally, it has to be fish- and plant-based, because you won ' t have \textbf{the} space or \textbf{the} resources for a dairy diet.
And finally, we are going to deal with all \textbf{the} waste locally, with compost, recycling, and turning \textbf{the} waste into energy.
So imagine where a traditional urban master plan, you typically draw \textbf{the} street grid where \textbf{the} cars can drive and \textbf{the} building plots where you can put some buildings.
This master plan, we sat down with a handful of scientists and basically started with all of \textbf{the} \textbf{renewable}, available natural resources, and then we started channeling \textbf{the} flow of resources through this kind of human- made ecosystem or this kind of urban metabolism.
So it ' s going to be modular, it ' s going to be buoyant, it ' s going to be designed to resist a tropical storm.
You can prefabricate it at scale, and tow it to dock with others, to form a small community.
We ' re designing these kind of coastal additions, so that even if it ' s modular and rational, each \textbf{island} can be unique with its own coastal landscape. \textbf{The} architecture has to remain relatively low to keep \textbf{the} center of gravity buoyant.
We ' re going to take all of \textbf{the} \textbf{agriculture} and use it to also create social space so you can actually \textbf{enjoy} \textbf{the} permaculture gardens.
We ' re designing it for \textbf{the} tropics, so all of \textbf{the} roofs are maximized to harvest solar power and to shade from \textbf{the} sun.
All \textbf{the} materials are going to be light and \textbf{renewable}, like bamboo and wood, which is also going to create this charming, warm environment.
And any architecture is supposed to be able to fit on this platform.
Underneath we have all \textbf{the} storage inside \textbf{the} pontoon, almost like a mega version of \textbf{the} student housings that we ' ve already worked with.
We have all \textbf{the} storage for \textbf{the} energy that ' s produced, all of \textbf{the} water storage and remediation.
We are sort of dealing with all of \textbf{the} waste and \textbf{the} composting.
And we also have some backup farming with aeroponics and hydroponics.
So imagine almost like a vertical section through this landscape that goes from \textbf{the} air above, where we have vertical farms ; below, we have \textbf{the} aeroponics and \textbf{the} aquaponics.
Even further below, we have \textbf{the} ocean farms and where we tie \textbf{the} \textbf{island} to \textbf{the} ground, we ' re using biorock to create new reefs to regenerate \textbf{habitat}.
So think of this small \textbf{island} for 300 people.
It can then group together to form a cluster or a neighborhood that then can sort of group together to form an entire city for 10, 000 people.
And you can imagine if this floating city flourishes, it can sort of grow like a culture in a petri dish.
So one of \textbf{the} first places we are looking at placing this, or anchoring this floating city, is in \textbf{the} Pearl River delta.
So imagine this kind of canopy of photovoltaics on this archipelago floating in \textbf{the} \textbf{sea}.
As you sail towards \textbf{the} \textbf{island}, you will see \textbf{the} maritime residents moving around on alternative forms of aquatic \textbf{transportation}.
You come into this kind of community port.
You can roam around in \textbf{the} permaculture gardens that are productive landscapes, but also social landscapes. \textbf{The} \textbf{greenhouses} also become orangeries for \textbf{the} cultural life of \textbf{the} city, and below, under \textbf{the} \textbf{sea}, it ' s teeming with life of farming and science and social spaces.
So in a way, you can imagine this community port is where people gather, both by day and by night.
And even if \textbf{the} first one is designed for \textbf{the} tropics, we also imagine that \textbf{the} architecture can adapt to any culture, so imagine, like, a Middle Eastern floating city or Southeast Asian floating city or maybe a Scandinavian floating city one day.
So maybe just to \textbf{conclude}. \textbf{The} human body is 70 percent water.
And \textbf{the} \textbf{surface} of our \textbf{planet} is 70 percent water.
And it ' s rising.
And even if \textbf{the} whole world woke up tomorrow and became carbon-neutral over night, there are still \textbf{island} nations that are destined to sink in \textbf{the} \textbf{seas}, unless we also develop alternate forms of floating human \textbf{habitats}.
And \textbf{the} only constant in \textbf{the} \textbf{universe} is change.
Our world is always changing, and right now, our climate is changing.
No matter how critical \textbf{the} crisis is, and it is, this is also our collective human superpower.
That we have \textbf{the} power to adapt to change and we have \textbf{the} power to give form to our future.

% matched lemmas: agriculture, conclude, consumption, enjoy, greenhouse, habitat, highway, island, organic, planet, pocket, renewable, sea, slide, slope, stack, surface, the, ticket, transportation, universe, wet
\end{textsample}
